\subsection{Special Tree Structures}
\label{sec:special_structures}

This section examines the worst-case performance of various special tree structures. All analyses focus on structural properties and decision tree depth, with uniform query costs.

\subsubsection{Star Graphs}
\label{sec:star_graphs}

\begin{table}[H]
\centering
\caption{Star graphs (worst case depth)}
\label{tab:star_uniform}
\begin{tabular}{|c|c|c|c|c|}
\hline
$n$ & Avg Depth & Max Depth & Avg Time (ms) & Std Dev \\
\hline
% Data from: star_uniform.csv
\csvreader[head to column names]{data/star_uniform.csv}{}{%
\n & \avgdepth & \maxdepth & \avgtime & \stddev \\
}
\hline
\end{tabular}
\end{table}
\begin{figure}[H]
\centering
\begin{tikzpicture}
\begin{axis}[
    xlabel={Tree size ($n$)},
    ylabel={Average depth},
    width=0.48\textwidth,
    height=6cm,
    grid=major,
    legend pos=north west,
    mark size=2pt
]
\addplot[blue, mark=*, thick] table[x=n, y=avgdepth, col sep=comma] {data/star_uniform.csv};
\legend{Avg Depth}
\end{axis}
\end{tikzpicture}
\hfill
\begin{tikzpicture}
\begin{axis}[
    xlabel={Tree size ($n$)},
    ylabel={Runtime (ms)},
    width=0.48\textwidth,
    height=6cm,
    grid=major,
    legend pos=north west,
    mark size=2pt
]
\addplot[red, mark=square*, thick] table[x=n, y=avgtime, col sep=comma] {data/star_uniform.csv};
\legend{Avg Time}
\end{axis}
\end{tikzpicture}
\caption{Performance metrics: Star graphs}
\label{fig:star_uniform}
\end{figure}

\textbf{Observations:} Star graphs achieve optimal depth of $\lceil \log_2 n \rceil$ due to their balanced structure with one central vertex.

\subsubsection{Caterpillar Trees}
\label{sec:caterpillar_trees}

\begin{table}[H]
\centering
\caption{Caterpillar trees (worst case depth)}
\label{tab:caterpillar_uniform}
\begin{tabular}{|c|c|c|c|c|}
\hline
$n$ & Avg Depth & Max Depth & Avg Time (ms) & Std Dev \\
\hline
% Data from: caterpillar_uniform.csv
\csvreader[head to column names]{data/caterpillar_uniform.csv}{}{%
\n & \avgdepth & \maxdepth & \avgtime & \stddev \\
}
\hline
\end{tabular}
\end{table}
\begin{figure}[H]
\centering
\begin{tikzpicture}
\begin{axis}[
    xlabel={Tree size ($n$)},
    ylabel={Average depth},
    width=0.48\textwidth,
    height=6cm,
    grid=major,
    legend pos=north west,
    mark size=2pt
]
\addplot[blue, mark=*, thick] table[x=n, y=avgdepth, col sep=comma] {data/caterpillar_uniform.csv};
\legend{Avg Depth}
\end{axis}
\end{tikzpicture}
\hfill
\begin{tikzpicture}
\begin{axis}[
    xlabel={Tree size ($n$)},
    ylabel={Runtime (ms)},
    width=0.48\textwidth,
    height=6cm,
    grid=major,
    legend pos=north west,
    mark size=2pt
]
\addplot[red, mark=square*, thick] table[x=n, y=avgtime, col sep=comma] {data/caterpillar_uniform.csv};
\legend{Avg Time}
\end{axis}
\end{tikzpicture}
\caption{Performance metrics: Caterpillar trees}
\label{fig:caterpillar_uniform}
\end{figure}


\subsubsection{Balanced Binary Trees}

\begin{table}[H]
\centering
\caption{Balanced binary trees (worst case depth)}
\label{tab:balanced_uniform}
\begin{tabular}{|c|c|c|c|c|}
\hline
$n$ & Avg Depth & Max Depth & Avg Time (ms) & Std Dev \\
\hline
% Data from: balanced_uniform.csv
\csvreader[head to column names]{data/balanced_uniform.csv}{}{%
\n & \avgdepth & \maxdepth & \avgtime & \stddev \\
}
\hline
\end{tabular}
\end{table}
\begin{figure}[H]
\centering
\begin{tikzpicture}
\begin{axis}[
    xlabel={Tree size ($n$)},
    ylabel={Average depth},
    width=0.48\textwidth,
    height=6cm,
    grid=major,
    legend pos=north west,
    mark size=2pt
]
\addplot[blue, mark=*, thick] table[x=n, y=avgdepth, col sep=comma] {data/balanced_uniform.csv};
\legend{Avg Depth}
\end{axis}
\end{tikzpicture}
\hfill
\begin{tikzpicture}
\begin{axis}[
    xlabel={Tree size ($n$)},
    ylabel={Runtime (ms)},
    width=0.48\textwidth,
    height=6cm,
    grid=major,
    legend pos=north west,
    mark size=2pt
]
\addplot[red, mark=square*, thick] table[x=n, y=avgtime, col sep=comma] {data/balanced_uniform.csv};
\legend{Avg Time}
\end{axis}
\end{tikzpicture}
\caption{Performance metrics: Balanced binary trees}
\label{fig:balanced_uniform}
\end{figure}


\subsubsection{Spider Graphs}

\begin{table}[H]
\centering
\caption{Spider graphs (worst case depth)}
\label{tab:spider_uniform}
\begin{tabular}{|c|c|c|c|c|}
\hline
$n$ & Avg Depth & Max Depth & Avg Time (ms) & Std Dev \\
\hline
% Data from: spider_uniform.csv
\csvreader[head to column names]{data/spider_uniform.csv}{}{%
\n & \avgdepth & \maxdepth & \avgtime & \stddev \\
}
\hline
\end{tabular}
\end{table}
\begin{figure}[H]
\centering
\begin{tikzpicture}
\begin{axis}[
    xlabel={Tree size ($n$)},
    ylabel={Average depth},
    width=0.48\textwidth,
    height=6cm,
    grid=major,
    legend pos=north west,
    mark size=2pt
]
\addplot[blue, mark=*, thick] table[x=n, y=avgdepth, col sep=comma] {data/spider_uniform.csv};
\legend{Avg Depth}
\end{axis}
\end{tikzpicture}
\hfill
\begin{tikzpicture}
\begin{axis}[
    xlabel={Tree size ($n$)},
    ylabel={Runtime (ms)},
    width=0.48\textwidth,
    height=6cm,
    grid=major,
    legend pos=north west,
    mark size=2pt
]
\addplot[red, mark=square*, thick] table[x=n, y=avgtime, col sep=comma] {data/spider_uniform.csv};
\legend{Avg Time}
\end{axis}
\end{tikzpicture}
\caption{Performance metrics: Spider graphs}
\label{fig:spider_uniform}
\end{figure}


\subsubsection{Broom Trees}

\begin{table}[H]
\centering
\caption{Broom trees (worst case depth)}
\label{tab:broom_uniform}
\begin{tabular}{|c|c|c|c|c|}
\hline
$n$ & Avg Depth & Max Depth & Avg Time (ms) & Std Dev \\
\hline
% Data from: broom_uniform.csv
\csvreader[head to column names]{data/broom_uniform.csv}{}{%
\n & \avgdepth & \maxdepth & \avgtime & \stddev \\
}
\hline
\end{tabular}
\end{table}
\begin{figure}[H]
\centering
\begin{tikzpicture}
\begin{axis}[
    xlabel={Tree size ($n$)},
    ylabel={Average depth},
    width=0.48\textwidth,
    height=6cm,
    grid=major,
    legend pos=north west,
    mark size=2pt
]
\addplot[blue, mark=*, thick] table[x=n, y=avgdepth, col sep=comma] {data/broom_uniform.csv};
\legend{Avg Depth}
\end{axis}
\end{tikzpicture}
\hfill
\begin{tikzpicture}
\begin{axis}[
    xlabel={Tree size ($n$)},
    ylabel={Runtime (ms)},
    width=0.48\textwidth,
    height=6cm,
    grid=major,
    legend pos=north west,
    mark size=2pt
]
\addplot[red, mark=square*, thick] table[x=n, y=avgtime, col sep=comma] {data/broom_uniform.csv};
\legend{Avg Time}
\end{axis}
\end{tikzpicture}
\caption{Performance metrics: Broom trees}
\label{fig:broom_uniform}
\end{figure}


\subsubsection{Lobster Trees}

\begin{table}[H]
\centering
\caption{Lobster trees (worst case depth)}
\label{tab:lobster_uniform}
\begin{tabular}{|c|c|c|c|c|}
\hline
$n$ & Avg Depth & Max Depth & Avg Time (ms) & Std Dev \\
\hline
% Data from: lobster_uniform.csv
\csvreader[head to column names]{data/lobster_uniform.csv}{}{%
\n & \avgdepth & \maxdepth & \avgtime & \stddev \\
}
\hline
\end{tabular}
\end{table}
\begin{figure}[H]
\centering
\begin{tikzpicture}
\begin{axis}[
    xlabel={Tree size ($n$)},
    ylabel={Average depth},
    width=0.48\textwidth,
    height=6cm,
    grid=major,
    legend pos=north west,
    mark size=2pt
]
\addplot[blue, mark=*, thick] table[x=n, y=avgdepth, col sep=comma] {data/lobster_uniform.csv};
\legend{Avg Depth}
\end{axis}
\end{tikzpicture}
\hfill
\begin{tikzpicture}
\begin{axis}[
    xlabel={Tree size ($n$)},
    ylabel={Runtime (ms)},
    width=0.48\textwidth,
    height=6cm,
    grid=major,
    legend pos=north west,
    mark size=2pt
]
\addplot[red, mark=square*, thick] table[x=n, y=avgtime, col sep=comma] {data/lobster_uniform.csv};
\legend{Avg Time}
\end{axis}
\end{tikzpicture}
\caption{Performance metrics: Lobster trees}
\label{fig:lobster_uniform}
\end{figure}


\subsubsection{Full m-ary Trees}
\label{sec:special_mary}

\begin{table}[H]
\centering
\caption{Full m-ary trees ($m=3$, worst case depth)}
\label{tab:mary_uniform}
\begin{tabular}{|c|c|c|c|c|}
\hline
$n$ & Avg Depth & Max Depth & Avg Time (ms) & Std Dev \\
\hline
% Data from: mary_uniform.csv
\csvreader[head to column names]{data/mary_uniform.csv}{}{%
\n & \avgdepth & \maxdepth & \avgtime & \stddev \\
}
\hline
\end{tabular}
\end{table}
\begin{figure}[H]
\centering
\begin{tikzpicture}
\begin{axis}[
    xlabel={Tree size ($n$)},
    ylabel={Average depth},
    width=0.48\textwidth,
    height=6cm,
    grid=major,
    legend pos=north west,
    mark size=2pt
]
\addplot[blue, mark=*, thick] table[x=n, y=avgdepth, col sep=comma] {data/mary_uniform.csv};
\legend{Avg Depth}
\end{axis}
\end{tikzpicture}
\hfill
\begin{tikzpicture}
\begin{axis}[
    xlabel={Tree size ($n$)},
    ylabel={Runtime (ms)},
    width=0.48\textwidth,
    height=6cm,
    grid=major,
    legend pos=north west,
    mark size=2pt
]
\addplot[red, mark=square*, thick] table[x=n, y=avgtime, col sep=comma] {data/mary_uniform.csv};
\legend{Avg Time}
\end{axis}
\end{tikzpicture}
\caption{Performance metrics: Full m-ary trees}
\label{fig:mary_uniform}
\end{figure}

Full m-ary trees have all internal nodes with degree $m$ and provide a balanced structure.

\textbf{Observations:}
\begin{itemize}
    \item Different tree structures exhibit varying worst-case search complexities
    \item Balanced structures (star, spider, balanced binary) achieve near-optimal depth
    \item Path-like structures (caterpillar, broom) have higher maximum depth
    \item Structural topology dominates worst-case performance
    \item Algorithm efficiently computes decision trees across diverse topologies
\end{itemize}
